\section{The Case for the Father of Latin Theology}

``Father of Latin theology'' is a received title, not an ancient one. This
section weighs it as a historical claim: what Latin Christianity existed
before him, what he actually coined, what he owed to the Greeks, and where the
title's honest limits lie.

\subsection{What was there before him}

Latin-speaking Christianity did not begin with Tertullian. The Scillitan
martyrs of 180 answered a Latin interrogation and carried ``books and epistles
of Paul'' in their case (\work{Acts of the Scillitan Martyrs}); Tertullian's
own biblical quotations presuppose existing Latin renderings of Scripture
circulating in Africa---he is, in Chapman's phrase, ``the first witness to the
existence of a Latin Bible,'' witness, not maker. What did not exist, so far
as any surviving text shows, was Latin theology: argument, in Latin, about
God, Christ, church, and sacrament, conducted at full intellectual pressure.
The one possible rival is the elegant dialogue \work{Octavius} of Minucius
Felix, another African; whether it precedes or follows the \work{Apology} is
a genuinely unresolved scholarly question, and this study does not adjudicate
it. But the \work{Octavius} is a philosophical apology that barely names
Christian doctrine; even granting it priority, the theological corpus in
Latin begins with Tertullian.

\subsection{The coinages}

The hard core of the title is lexical, and the specimens are on his pages at
verified loci. \textit{Trinitas} as the name of the three-personed God does
its first surviving systematic work in \work{Against Praxeas} 2--3
(\textit{quae unitatem in trinitatem disponit}); \textit{substantia} and
\textit{persona}, put asymmetrically to God's oneness and threeness in the
same chapter (\textit{tres\ldots\ unius autem substantiae}), became and
remain the Latin dogmatic pair; \textit{una persona} with \textit{duplex
status} for Christ (\work{Against Praxeas} 27) gave the West its
Christological syntax; \textit{sacramentum} was bent from the soldier's oath
to the church's mysteries. Around the technical stock stand the coined
sentences that Latin Christianity never stopped repeating: \textit{semen est
sanguis Christianorum} (``the blood of Christians is seed,'' \work{Apology}
50); \textit{o testimonium animae naturaliter Christianae} (``O noble
testimony of the soul by nature Christian!'', \work{Apology} 17); \textit{caro
salutis est cardo} (``the flesh is the hinge of salvation,'' \work{On the
Resurrection of the Flesh} 8); \textit{certum est, quia impossibile}
(\work{On the Flesh of Christ} 5). Philologists have counted his neologisms
in the hundreds; the counts vary with the criteria and none was verified for
this study, so the claim is left qualitative: Chapman's older caution---``to
some extent, how great we cannot tell, he must have invented a theological
idiom and have coined new expressions''---remains the honest form of it.

\subsection{The debt to the Greeks}

None of this was creation from nothing. Tertullian read Greek, wrote Greek
works now lost, and worked with the Greek theology of his century all over
his desk. Theophilus of Antioch, a generation before him, had already called God, his
Word, and his Wisdom a ``Trinity'' in Greek---\textit{trias}: ``the three
days which were before the luminaries, are types of the Trinity, of God, and
His Word, and His wisdom'' (\work{To Autolycus} 2.15). Justin and the apologists had built the case from the soul and from
prophecy; Irenaeus had built the argument from apostolic succession that
underlies the \work{Prescription}---Chapman again: ``the whole luminous
argument is founded on the first chapters of St.~Iren\ae us's third book.''
Nor was the pagan inheritance disowned in practice, whatever
\work{Prescription} 7 thundered: in his psychology he can cite the Stoic
moralist warmly and possessively---``Just as Seneca says, whom we so often
find on our side'' (\work{A Treatise on the Soul} 20)---and his doctrines of
soul and substance are worked in Stoic materials throughout. What Tertullian
did was not to invent these arguments but to re-found them in
another language's categories: where Greek theology reached for philosophical
abstraction, he reached for the law-court and the army camp---person,
substance, prescription, oath. The choice of metaphors was itself a theology,
and it is the one the Latin West inherited.

\subsection{A style that taught Latin to argue about God}

The vocabulary travelled because the sentences carried it. Tertullian's
Latin is compressed, antithetical, and built for quotation: the paradox
sequence of \work{On the Flesh of Christ} 5, the balanced taunt of \work{On
Monogamy} 1 (``Heretics do away with marriages; Psychics accumulate them''),
the courtroom snap of \work{Apology} 2's ``O miserable deliverance,'' the
epigrams that this study has met on nearly every page. Antiquity itself
registered the cost of the compression: Lactantius, a professional
rhetorician, complained that in eloquence Tertullian ``had little readiness,
and was not sufficiently polished, and very obscure'' (\work{Divine
Institutes} 5.1)---and antiquity kept quoting him anyway, because obscurity
of texture and memorability of phrase turned out to be different things.
The judgment that matters historically is not aesthetic: his
manner---definition, distinction, antithesis, maxim---is already the working
method of Latin theology, and the first reader known to have studied him
daily was the next great African theologian, Cyprian, whose treatises rework
his themes (see the reception section). To a remarkable degree, Latin
theology argues the way its first practitioner wrote.

\subsection{The verdict, with limits}

The title, then, survives audit in a precise form. He is not the first Latin
Christian, nor the first Christian writer in Africa, nor the inventor of the
Latin Bible, nor the theological equal, in later eyes, of the tradition he
equipped; and one candid limit must stand beside every list of coinages---the
identical vocabulary served him, within a few years, to argue positions the
church refused. But as the first person known to have thought through the
whole Christian system in Latin, and to have left the language permanently
stocked with the words the dogmas would need, he has no competitor. Jerome's
ranking---``chief of the Latin writers''---and Benedict XVI's 2007
characterization of him as the beginner of Christian Latin literature and
theology are, on this point, reception and history in agreement.
